\chapter{结论与展望}\label{Chap:chap6}

\section{研究总结}\label{Sec6.1}

本文提出了Swift算法——一种启发式多信号融合自适应网络拥塞控制方法。"Swift"（雨燕）象征着算法追求的设计目标：如同雨燕般轻盈敏捷，实现高吞吐量（High Throughput）与低延迟（Low Latency）的平衡。Swift的窗口决策由确定性规则在每个确认事件上于协议栈之外求解，不依赖训练样本与神经网络推理，行为可逐分支解释、结果可完全复现。

本文的研究历程可以概括为以下几个阶段：

\textbf{问题分析阶段}：首先深入分析了远距离传输、蜂窝/无线/卫星等多接入链路以及手机、笔记本、台式机和边缘设备等异构终端并存场景中的拥塞控制挑战。在此基础上，系统性地分析了传统拥塞控制算法（基于丢包的和基于延迟的）的优缺点、ECN机制的潜力与利用不足，以及学习型拥塞控制方法在训练依赖、推理开销与可审计性方面的工程障碍，确定了以协议栈内部信号驱动的白盒启发式设计作为研究方向。

\textbf{方案设计阶段}：基于问题分析的结论，设计了Swift算法的核心组件。从状态设计入手，提出了11维有效观测子空间（嵌入于15元素状态传输容器，另含4项元数据通道）；然后设计了多信号融合拥塞检测方法，采用分层优先级机制区分超时、ECN驱动的窗口缩减与普通丢包；接着设计了差异化拥塞响应机制和启发式反馈自适应参数调优机制；最后将各组件整合为以$\alpha \times BDP$为目标的融合控制策略。

\textbf{实现与验证阶段}：在ns-3.40网络仿真平台上实现了Swift算法的原型系统，仿真进程与决策进程经ZeroMQ同步请求—应答信道交互；瓶颈网关配置启用ECN标记的RED主动队列管理，并对全部被比较协议对称开启ECN以保证对照公平性。实验矩阵覆盖近端高速、共享收敛链路、无线局域网接入、蜂窝移动接入、城域与长途/跨域广域网及LEO/GEO卫星通信共36个场景，在"纯TCP"与"TCP $+$ 平均负载为瓶颈32\%（峰值64\%、50\%占空比）的UDP Burst"两类设置下对四种协议各运行一次，共288次仿真。全部指标由FlowMonitor产物重新导出并仅在前向数据流上定义；288条记录全部通过完整性、物理相容性与取值域审计，无数据点被排除。

\textbf{分析与优化阶段}：对实验结果进行了深入分析。Swift在全部36个纯TCP场景与36个UDP突发配对场景中均取得最高聚合吞吐量（相对最强基线约+2\%$\sim$+6\%），时延、丢包与公平性指标不劣于基线；同时如实识别了适用边界：吞吐增益为小幅一致优势而非数量级改善，个别深缓冲场景的时延高于延迟最优基线最高约10\%，丢包率与NewReno相当而非系统性更低。据此给出了以排队延迟感知的$\alpha$调度与多时间尺度BDP估计为代表的改进方向。

本文的主要贡献和研究成果总结如下：

\textbf{（1）完整的11维观测空间设计}

本文设计了11维有效观测子空间（嵌入于15元素状态传输容器，另含4项跨进程路由用的元数据通道），将显式拥塞通知（ECN）状态、拥塞控制状态机状态、拥塞事件类型等TCP协议栈内部信息纳入决策状态表示。11维观测空间涵盖窗口状态（ssthresh、cwnd）、传输指标（segSize、segmentsAcked、bytesInFlight）、延迟测量（lastRtt、minRtt）、回调类型（callbackType）和拥塞控制状态（caState、caEvent、ecnState）等多个维度，为决策模块提供了完整的协议栈状态感知能力。

与现有方案通常采用的4--6维简化状态空间相比，Swift的11维观测空间能够提供更丰富、更精细的网络状态信息，为精准的拥塞控制决策奠定了基础。

\textbf{（2）多信号融合拥塞检测与差异化响应机制}

本文提出了一种多信号融合的拥塞检测方法，采用两级优先级机制综合分析多种拥塞信号。第一优先级作用于协议栈的显式降窗回调（callbackType=0）内部，按拥塞状态机状态与ECN状态完成三路语义判别：caState=CA\_LOSS判定为重传超时（保留因子0.50）；ecnState为ECN\_CE\_RCVD或ECN\_ECE\_RCVD、或caState=CA\_CWR判定为ECN回显驱动的窗口缩减（保留因子0.75）；其余情形判定为快速重传丢包（保留因子0.70）。第二优先级在正常确认路径上巡检ECN状态。该三路判别是必要的：该回调本身并不携带降窗原因，若笼统映射为"丢包"，ECN提供的早期预警在响应环节会被当作队列已溢出处理，从而抵消ECN机制的时序优势。同时引入连续缩减保护机制，限制最多3次连续窗口缩减后保持窗口不变，并在显式降窗后冻结4个ACK事件中的窗口增长；连续缩减计数器仅在窗口真实增长时清零，以免冻结期内的每个ACK事件反复复位保护计数。

为避免对瞬态扰动的过度响应，Swift仅在上述两级判别命中时执行拥塞响应：RTT膨胀、CA\_DISORDER与CA\_RECOVERY状态不作为独立的降窗触发条件，其中RTT膨胀转而用于$\alpha$自适应与慢启动退出判定。

\textbf{（3）启发式反馈自适应的融合控制策略}

本文实现了启发式反馈自适应的参数调优机制，根据三类信号动态调整乘性增加因子$\alpha$：其一，根据最小RTT构造路径感知的RTT膨胀动态阈值，使长RTT路径不因正常传播抖动而过早收紧窗口；其二，以性能反馈值快线EMA（$\eta_f=0.15$）相对慢线基线EMA（$\eta_b=0.02$）的偏离量判断近期控制效果，容差$m=\max(0.25, 0.1|R_{base}|)$随基线幅值自适应缩放——采用相对判据而非固定绝对阈值是必要的，因为环境侧反馈值在绝大多数确认事件上为正，绝对阈值会使调节退化为持续推向上界的单向棘轮；其三，在连续增长超过8次时适度增大$\alpha$以提升带宽利用率。$\alpha$的取值被限制在$[0.85, 1.30]$范围内，基准值为1.10，避免在无线、蜂窝和广域网场景中持续堆积队列。

融合控制策略将基于带宽延迟积（BDP）估计的目标窗口与当前拥塞窗口状态结合。BDP采用两级结构估计：第一级按时间窗口差分计算投递速率样本，窗口跨度不小于两倍最小RTT并被限制在5\,ms至1\,s之间，使分子的确认字节数与分母的时间区间自洽，消除了旧版逐ACK公式在长RTT路径上的系统性低估；第二级将样本压入容量40的滑动窗口并取最大值作为瓶颈带宽，再乘以最小RTT得到稳健估计。在慢启动阶段，目标窗口设为BDP的2倍，并通过HyStart风格RTT膨胀检测退出慢启动；在拥塞避免阶段，算法围绕$\alpha \times BDP$目标窗口进行有界步长的渐进拉升或回落。窗口上界采用$\max(4 \times BDP, 200 \times MSS)$，下界为$4 \times MSS$。每流独立状态管理用于维护历史度量缓冲区、拥塞事件计数、自适应参数和降窗后冻结计数等运行状态，为多流场景中的差异化控制提供实现基础。

在实验验证方面，本文在ns-3.40网络仿真平台上实现了Swift算法，并在36个场景、288次仿真（两种设置各144次）的全矩阵下进行了系统性对比实验。实验结果表明：

\begin{itemize}
    \item 纯TCP设置下Swift在\textbf{全部36个场景}取得最高聚合吞吐量：相对最强基线$+2.1\%\sim+5.8\%$（均值$+4.1\%$），相对NewReno/CUBIC均值$+6.0\%$；增益在微秒级近端高速、无线/蜂窝接入与远距离（含LEO/GEO卫星、长途与跨域WAN）链路三类路径上均匀分布。
    \item 平均单向时延与基线同量级：相对NewReno/CUBIC均值平均低约2.3\%（36个场景中24个更低、无一超过+1.7\%），相对延迟最低的BBR平均高约0.8\%（个别深缓冲场景最高约10.4\%）。
    \item 四协议丢包率全部不超过0.026\%（UDP突发下不超过0.045\%）；Swift平均丢包率与NewReno相当、略低于CUBIC/BBR，本文如实说明其并不系统性低于全部基线。
    \item Swift的Jain公平性指数均值0.9989（纯TCP）/0.9990（UDP突发）、最小0.9967/0.9940，均为四协议最高。
    \item 在GEO卫星场景（S19，基线单向时延320\,ms）中，Swift纯TCP吞吐量8.78\,Mbps（相对NewReno提升7.5\%）；UDP突发下5.65\,Mbps（相对NewReno提升10.1\%），为四协议最高。LEO卫星与长途WAN场景同样领先。
    \item 在平均负载为瓶颈32\%的UDP突发下，四协议吞吐一致回落约31\%--32\%，无任何协议发生塌缩；Swift的绝对吞吐量在36/36个配对场景保持最高（$+2.2\%\sim+5.9\%$），时延膨胀（+2.9\%）与新增丢包（+0.013个百分点）与基线无系统分化。
    \item 旧版逐ACK投递速率估计器在长RTT路径上的窗口塌缩问题（历史数据中跨域WAN回落约29\%）已由时间窗口估计器修复，并在本次全矩阵数据中得到验证：跨域WAN（874.6\,Mbps）、长途WAN（881.2\,Mbps）等长RTT场景均为四协议最高。
\end{itemize}

\section{研究展望}\label{Sec6.2}

尽管Swift算法在36场景全矩阵下取得了方向一致的性能优势，但当前结论仍受限于ns-3仿真环境、同步决策路径与RED/ECN配置范围。以下内容作为后续研究方向，不作为本文已经完成的实现或实验结论。

\subsection{参数自适应与窗口控制的细化}

实验记录了两处可继续优化的边界：（1）Swift在个别深缓冲/低速率场景的相对时延抬升最高约10\%；（2）长RTT路径上的吞吐增益（+2.2\%$\sim$+4.8\%）低于近端场景。可探索的方向包括：

\begin{itemize}
    \item 排队延迟显式估计：引入$q_{est} = RTT_{smoothed} - RTT_{min}$作为$\alpha$调整的额外输入，使$\alpha$在排队延迟持续增长时自动收敛至1.0以下。
    \item 多时间尺度BDP估计：区分传播延迟和排队延迟的贡献，避免将排队延迟计入BDP估计。
    \item 受约束的带宽探测：在不突破当前$\alpha\in[0.85,1.30]$稳定性边界的前提下，为长RTT路径设计更细粒度的窗口拉升规则。
    \item 机制消融实验：分别关闭ECN/丢包差异化保留因子、连续缩减保护与降窗冻结，量化各机制的独立贡献；当前数据无法完成该分离。
\end{itemize}

\subsection{策略表达能力扩展}

当前Swift算法采用基于规则的启发式决策机制，具有可解释性和低延迟优势，但表示能力受人工规则结构限制。后续研究可在不改变本文实验结论的前提下，探索离线强化学习训练神经网络策略、再通过模型蒸馏提取为紧凑规则的路线；也可引入元学习方法加快不同网络环境之间的参数迁移。相关方法需要在完整实验矩阵、异常过滤和可复现日志基础上重新评估，在获得结果之前不作为本文系统的组成部分。

\subsection{多目标优化扩展}

当前Swift的优化目标主要是高吞吐量、低延迟和低丢包。在移动终端、边缘计算和多业务并发场景中，后续研究可进一步引入公平性、能耗、可靠性和业务优先级等目标。多目标优化存在吞吐、延迟、公平性和能耗之间的冲突，需要基于约束优化方法重新设计反馈函数与评价体系。

\subsection{真实网络部署与验证}

本文实验主要在ns-3仿真平台上进行，与真实网络环境存在差距。后续部署验证需要将算法实现为Linux内核TCP模块或QUIC用户态协议栈，并在覆盖近端高速、蜂窝/无线接入、广域网与卫星链路的真实或半实物测试床上重新测量。真实环境中的NIC offload、中断合并、NUMA效应、操作系统调度、无线链路波动和终端算力差异均可能影响实际性能，因此不能直接将本文仿真结果外推为生产部署效果。

\subsection{实验方法的扩展}

受单次确定性运行限制，本文未报告跨种子置信区间。后续应使用\texttt{main.py --num-seeds}对关键场景做多种子重复并给出均值与最坏值；同时可扩展队列管理策略（丢尾、CoDel）与突发强度/占空比扫描，检验结论对配置的稳健性。

\subsection{安全性、鲁棒性与跨层协同}

后续研究需要进一步验证算法在异常输入、通信中断、状态损坏和恶意跨流量下的行为边界。可行方向包括：输入合法性检查、异常状态回退、保守策略兜底、形式化安全性质描述，以及与应用层、链路层和操作系统调度器的协同机制。上述方向均需要独立实现和实验验证，本文不将其作为当前系统已经具备的能力。

\subsection{异构网络与新兴传输协议适配}

远距离传输和异构终端场景通常跨越有线、无线、蜂窝、卫星和边缘接入等多类链路。后续研究可进一步考察无线非拥塞丢包识别、卫星链路长RTT适配、5G突发抖动处理、终端算力约束下的轻量化执行，以及将多信号融合、差异化响应和自适应调参思想迁移到QUIC用户态协议栈中的可行性。相关扩展应基于新的实现版本、完整实验矩阵和异常过滤流程重新评估。

\section{应用场景展望}\label{Sec6.3}

基于本文仿真结果，Swift更适合作为面向远距离传输、异构终端和多接入链路的拥塞控制研究原型，而不是已经完成真实网络生产部署验证的协议。其潜在应用场景包括：

\begin{itemize}
    \item \textbf{跨地域内容传输}：在长RTT路径上维持更高吞吐并控制队列驻留（跨域/长途WAN与GEO/LEO场景实测领先）。
    \item \textbf{蜂窝与无线接入}：利用ECN、RTT和在途字节等多源信号提高对接入链路波动的适应能力。
    \item \textbf{移动终端与边缘应用}：面向手机、笔记本、桌面终端和边缘设备并存的场景，探索更稳健的参数自适应机制。
    \item \textbf{卫星与广域链路}：在高传播时延和高BDP路径上保持带宽估计有效并减少拥塞恢复代价。
\end{itemize}

上述应用方向体现了算法设计的潜在价值，但不意味着本文仿真结果可以直接等同于生产网络收益。实际部署前仍需要完成内核或用户态协议栈实现、真实测试床验证、跨流公平性评估和长期稳定性评估。

\section{研究意义与社会价值}\label{Sec6.4}

本研究的理论意义主要体现在三个方面：首先，提出了面向远距离和异构终端传输的11维有效观测子空间，为协议栈信号驱动的拥塞控制状态设计提供了参考；其次，提出了多信号融合拥塞检测与差异化响应机制，为区分ECN、丢包和超时等不同拥塞信号提供了可解释规则；最后，提出了结合BDP估计、路径感知$\alpha$调节、降窗后冻结和连续缩减保护的白盒启发式融合控制框架，验证了"确定性规则+轻量反馈自适应"设计点在跨场景一致性上的可行性。

本研究的实际价值在于：通过ns-3.40构建了一个可复现实验原型，展示了多信号融合与启发式反馈自适应在长距离、多接入和异构终端传输中的潜力——在RED/ECN对称配置下，该方法在36个纯TCP场景与36个突发配对场景中均取得最高吞吐量，且不伴随时延、丢包或公平性的系统性劣化。若后续能够在真实协议栈和真实网络环境中完成验证，该类方法有望提升复杂网络传输中的吞吐、延迟、丢包和稳定性综合表现。
